%%% FIELD proof-tight-moment-radius-forces-zero-means
Put \(q_0=4+\delta\), and suppose
\[
M_{\mathrm{mom}}=(2v_\star)^{q_0/2}.
\]
If \(\mathcal M=\varnothing\), the asserted universal conclusion is vacuous. Otherwise fix arbitrary \(P\in\mathcal M\) and \(c\in[C]\). By \(\delta>0\) and \(v_\star>0\), one has \(q_0>0\) and \(2v_\star>0\). Hence the definition \(\operatorname{def:effective\text{-}loading\text{-}mean\text{-}radius}\) gives
\[
\begin{aligned}
b_\lambda
&=\left(M_{\mathrm{mom}}^{2/q_0}-2v_\star\right)_+^{1/2}\\
&=\left(\left\{(2v_\star)^{q_0/2}\right\}^{2/q_0}-2v_\star\right)_+^{1/2}\\
&=\left(2v_\star-2v_\star\right)_+^{1/2}=0.
\end{aligned}
\]
The zero-radius boundary conclusion in \(\operatorname{lem:effective\text{-}loading\text{-}mean\text{-}radius\text{-}bound}\), applied to this \(P\) and \(c\), therefore yields
\[
H_c=v_\star I_2,
\qquad
m_c=0,
\qquad
\varepsilon_i=0\quad\text{almost surely conditional on }C_i=c,
\qquad
\mu_c(P)=0.
\]
Because \(P\in\mathcal M\) and \(c\in[C]\) were arbitrary, \(\mu(P)=0\) for every \(P\in\mathcal M\). No step uses \(v^{\star}=v_\star\), so no equality between the covariance upper and lower radii is required.
